\documentclass[../../../../SoftwareRequirementsSpecification.tex]{subfiles}
\begin{document}
% Richiede le macro definite in src/macros/useCaseMacros.tex
% (incluse nel preambolo di SoftwareRequirementsSpecification.tex)

\ucstart
\begin{tabularx}{\textwidth}{|l|X|}
  \hline
  \ucid{PT-01} \\ \hline
  \ucname{Registrazione PT} \\ \hline
  \ucactor{PT} \\ \hline
  \ucdescription{Un personal trainer crea autonomamente il proprio
    account per poter utilizzare l'applicativo.} \\ \hline
  \ucprecond{L'utente non ha un account nel sistema.} \\ \hline
  \ucpostcond{Il PT ha un account nell'applicativo e può accedervi
    tramite il caso d'uso U-01. Non ha ancora alcun atleta
    associato.} \\ \hline
  \ucmainflow{
    \item L'utente atterra sulla pagina di registrazione.
    \item Il sistema mostra il form di inserimento dati (nome,
      cognome e mail).
    \item L'utente inserisce i dati.
    \item L'utente preme il pulsante "Registrati".
    \item Il sistema verifica che i dati siano validi e che la mail
      non sia già presente nel database.
    \item Il sistema crea l'account del PT.
    \item Il sistema mostra una pagina di conferma e indirizza
      l'utente al caso d'uso U-01 (Login).
  } \\ \hline
  \ucaltflow{
    \textbf{5a. Email già registrata:} \newline
    - Il sistema mostra un messaggio di errore ("Email già
    registrata"). \newline
    - Il flusso riprende dal punto 2. \newline
    \textbf{5b. Dati non validi:} \newline
    - Il sistema mostra un messaggio di errore ("Controllare i dati
    inseriti"), indicando i campi non validi (email in formato non
    valido, nome o cognome vuoti). \newline
    - Il flusso riprende dal punto 2.
  } \\ \hline
  \ucnote{
    La registrazione non stabilisce alcuna credenziale: il sistema è
    passwordless e l'unico meccanismo di accesso è il codice
    temporaneo di U-01. Per questo motivo il caso d'uso non prevede
    un passo di verifica dell'indirizzo di posta elettronica: il
    possesso della mail è dimostrato al primo accesso, dato che senza
    ricevere il codice temporaneo l'account resta inutilizzabile. Per
    lo stesso motivo il sistema non autentica l'utente al termine
    della registrazione, ma lo indirizza al login. \newline
    Il messaggio del flusso 5a non distingue fra una mail già usata
    da un altro PT e una mail di un atleta già registrato dal proprio
    PT: l'unicità di \texttt{email} su \texttt{User} vale sull'intera
    gerarchia, quindi uno stesso utente non può ricoprire entrambi i
    ruoli.
  } \\ \hline
\end{tabularx}
\end{document}
